\section{Structural exact algorithms}\label{s4:structural}

The hardness results of \cref{s3:hardness} do not make every large
pooling network intractable. Two decompositions give exact algorithms.
The first keeps a fixed number of nonlinear variables and aggregate
rows and partitions all other flows into independent polyhedra of fixed
dimension (\cref{s4:core-model,s4:pooling-basic,s4:bypass-structure}).
The second eliminates long bypass paths and keeps a fixed number of
actual boundary flows
(\cref{s4:planar,s4:attachments,s4:bounded-support}). Both allow
arbitrarily many variables; their hypotheses control different sources
of coupling. Throughout this section pooling means the standard three-layer model of
Section~\ref{s1:model}, with finite rational flow bounds, arbitrary
rational lower and upper throughput and attribute specifications, and
no pool--pool arcs. Unless restricted explicitly, objectives are rational
linear costs on actual arcs. All algorithmic conclusions concern binary
input size and exact algebraic output as in Section~\ref{s1:encoding}.

The two decompositions retain different information, and this
determines the objective scope of each result. The block construction
retains both aggregate masses and the objective value, and therefore
permits costs on every arc. The path construction retains only boundary
feasibility; it permits arbitrary costs only when exact equalities
reduce them to the retained coordinates.

\subsection{A fixed nonlinear core and independent polyhedral blocks}
\label{s4:core-model}

Fix integers $r,k\ge0$ and $d,D\ge1$. Let $C\subset\R^r$ be a compact
set specified by a quantifier-free rational polynomial formula of degree
at most $D$. For $j=1,\ldots,N$, let
\[
 P_j(x)=\{z\in\R^{d_j}:B_j(x)z\le b_j(x)\},\qquad 1\le d_j\le d.
\]
Every coordinate has an explicit finite rational lower and upper box
bound among these inequalities. Empty blocks and lower-dimensional
blocks are allowed. All entries of $B_j,b_j,A_j,b,c_j,c_0$ below are
rational polynomials in $x$ of degree at most $D$, with
$A_j(x)\in\R^{k\times d_j}$ and $b(x)\in\R^k$. Consider
\begin{equation}
 \begin{aligned}
 \min\quad &c_0(x)+\sum_{j=1}^N c_j(x)^Tz_j,\\
 \text{subject to}\quad &x\in C,\quad z_j\in P_j(x)\quad(1\le j\le N),\\
 &\sum_{j=1}^N A_j(x)z_j=b(x).
 \end{aligned}\label{s4:block-model}
\end{equation}
The row counts of the local systems and the number of blocks are part
of the input and may grow.

\begin{theorem}[Fixed core, local dimension, and aggregate dimension]
\label{s4:fixed-core}
For every fixed $(r,d,k,D)$, feasibility and global optimization of
\eqref{s4:block-model} are solvable in polynomial bit time. If feasible,
the algorithm returns an exact optimum and optimizer of polynomial
encoding length. All optimizer coordinates can be represented in one
real-algebraic extension of $\Q$ of polynomial degree. A fixed number
of aggregate inequalities can also be included.
\end{theorem}

The geometric ingredients are established: support directions select
compatible faces of Minkowski summands, and sums of polytopes in fixed
dimension have polynomial complexity
\citep{s4:gritzmann1993,s4:fukuda2004}. These results concern specified
polytopes. Here every summand changes with a nonlinear parameter, so
computing a Minkowski sum separately at each parameter would not give
a finite global algorithm. The additional construction is a polynomial
list of support formulas valid uniformly over the whole core, including
parameters where a local basis becomes singular or a block loses
dimension. The required sign enumeration, quantifier elimination, and
algebraic sampling are the algorithms of
\citet[Theorem~1.3.1 and Sections~3.1.3--3.2]
{basu1996-on-the-combinatorial-and-algebraic}.

\begin{proof}
If $N=0$, append a scalar block fixed to zero, with zero contribution
and cost. A block with no coordinates in an application is likewise
padded by one coordinate fixed to zero. We first prove the equation form.

\emph{Local vertices and their validity.}
Let $M_j$ be the number of rows in block $j$. For every $d_j$-row
subset $I$, form
\[
 \delta_{jI}(x)=\det B_{jI}(x),\qquad
 p_{jI}(x)=\operatorname{adj}(B_{jI}(x))b_{jI}(x).
\]
An identically zero determinant can be discarded. Otherwise the
candidate $v_{jI}=p_{jI}/\delta_{jI}$ is defined when
$\delta_{jI}\ne0$. For each row $a$, put
\begin{equation}
 e_{jIa}=B_{ja}p_{jI}-b_{ja}\delta_{jI}.
 \label{s4:vertex-feasibility}
\end{equation}
The candidate is feasible exactly when $\delta_{jI}\ne0$ and
$e_{jIa}/\delta_{jI}\le0$ for every $a$. These tests use signs of
polynomials of degree at most $(d+1)D$; the determinant and numerator
degrees are at most $dD$.

Every nonempty bounded polyhedron has a vertex. At any vertex its
active row normals span the ambient space: otherwise a nonzero
orthogonal direction, with sufficiently small positive and negative
steps, gives a feasible segment through that point. Thus some $d_j$
active rows are independent, even when the polyhedron has smaller
affine dimension or consists of one point. The candidates therefore
include every vertex, and a block is nonempty if and only if at least
one candidate passes its validity tests.

\emph{Support choices without a Cartesian product.}
Set $h=k+1$ and define
\[
 W_j(x)=\begin{pmatrix}A_j(x)\\c_j(x)^T\end{pmatrix},\quad
 T_j(x)=W_j(x)P_j(x),\quad
 w(x,v)=\begin{pmatrix}b(x)\\v-c_0(x)\end{pmatrix}.
\]
Here $v$ denotes an objective value. For a support direction
$u\in\R^h$, the score of a valid candidate is
\[
 \rho_{jI}(x,u)=\frac{n_{jI}(x,u)}{\delta_{jI}(x)},\qquad
 n_{jI}=u^TW_j(x)p_{jI}(x).
\]
Comparison of candidates $I,I'$ in the same block uses the determinant
signs and the sign of
\begin{equation}
 n_{jI}\delta_{jI'}-n_{jI'}\delta_{jI}.
 \label{s4:score-comparison}
\end{equation}
Let $\mathcal P$ contain these comparison polynomials, all determinants,
and all polynomials in \eqref{s4:vertex-feasibility}. Identically zero
polynomials have a remembered zero sign and need not be stored. Since
$\binom{M_j}{d_j}\le M_j^d$, the number of polynomials is polynomial
in the input size. Their degrees are bounded in terms of $d,D$: a
score comparison has degree at most $(2d+1)D+1$, counting the support
coordinates as variables. Their coefficient lengths are polynomial,
and they involve only the fixed
number $r+h$ of coordinates $(x,u)$.

Enumerate all realizable sign conditions of $\mathcal P$, including
zero signs, using fixed-dimensional real-algebraic sign enumeration.
There are polynomially many such conditions, and the enumeration has
polynomial bit complexity. A sign condition determines all valid
candidates and all score comparisons. If a block is empty, discard
the condition. Otherwise choose in every block the first maximizing
candidate in a fixed order. This choice is constant throughout the
realization of the sign condition; connectedness of that realization
is unnecessary. In particular, degeneracies and ties do not require
enumerating all tuples of local bases.

For a retained condition $\sigma$, abbreviate its selected determinants
and numerators to $\delta_j,n_j$ and put
\begin{align}
 H_\sigma(x)&=\prod_{j=1}^N\delta_j(x)^2,\\
 F_\sigma(x,v,u)&=u^Tw(x,v)H_\sigma(x)
 -\sum_{j=1}^N n_j(x,u)\delta_j(x)
                         \prod_{i\ne j}\delta_i(x)^2.
 \label{s4:cleared-support}
\end{align}
On this condition $H_\sigma>0$, so $F_\sigma\le0$ is equivalent to
\begin{equation}
 u^Tw(x,v)\le\sum_{j=1}^N
                   \max_{z\in P_j(x)}u^TW_j(x)z.
 \label{s4:support-inequality}
\end{equation}
Squared denominators preserve the inequality for either determinant
sign. The degree of $F_\sigma$ grows as $O(N)$, with constants depending
on the fixed parameters; it is not bounded independently of $N$.
Nevertheless it has polynomially many monomials because its number of
variables is fixed. Expanding the products takes polynomial bit time:
there are polynomially many multiplication steps and monomials, and
the coefficient lengths grow by at most the sum of the factor lengths
plus the logarithms of the numbers of terms accumulated. This also
bounds the coefficients of the sum in \eqref{s4:cleared-support}.

\emph{A formula in a fixed number of variables.}
Write $S_\sigma(x,u)$ for the conjunction of signs defining $\sigma$,
and set
\[
 \Phi(x,v,u)=\bigvee_{\sigma\text{ retained}}
           \bigl(S_\sigma(x,u)\wedge F_\sigma(x,v,u)\le0\bigr).
\]
We claim that the objective value $v$ is attainable at a specified
$x$ if and only if
\begin{equation}
 \forall u\in\R^h\quad\Phi(x,v,u).
 \label{s4:universal-support}
\end{equation}
If some block is empty, all conditions realized at $(x,u)$ are
discarded, for every $u$, so this formula is false. If all blocks are
nonempty, their image sum $T(x)=\sum_jT_j(x)$ is nonempty compact
convex. Its support function is the sum of their support functions.
A point belongs to a nonempty compact convex set exactly when it
satisfies every support inequality, by separating any point outside
the set. Thus \eqref{s4:universal-support} says precisely that
$w(x,v)\in T(x)$. This is equivalent to the existence of the block
vectors with the required aggregate and objective values, proving
the claim.

Let $E(x,v)$ abbreviate $x\in C$ and
\eqref{s4:universal-support}. The formula
\[
 \exists x\ E(x,v)
\]
defines the attainable objective values. Its total variable count is
$r+h+1$, and its degree, length, and rational coefficient encoding
are polynomial. Fixed-dimensional quantifier elimination therefore
produces a univariate formula in polynomial bit time, even with the
degree growth noted above. The original feasible set is closed in the
product of the compact core and the finite block boxes; hence it is
compact. The objective is continuous, so its nonempty value set is
compact and has a least element. Univariate root isolation decides
emptiness and finds that element exactly.

For joint optimizer recovery one can instead eliminate quantifiers in
\begin{equation}
 E(x,v)\ \wedge\
 \neg\exists x',v'\,[E(x',v')\wedge v'<v].
 \label{s4:optimal-sample}
\end{equation}
Rename the bound support variables in each occurrence of $E$.
The total number of variables remains fixed. Algebraic sampling gives
an optimal pair $(x^*,v^*)$ with polynomial-size rational univariate
representation: there is one real algebraic $\alpha$ of polynomial
degree, with isolating data, and rational functions of $\alpha$
representing all these coordinates. Equivalently, they lie in the
common represented field $F=\Q(\alpha)$.

\emph{Recovering all blocks.}
Fix $x=x^*$ in the local inequalities and aggregate equations, and
append
\[
 \sum_jc_j(x^*)^Tz_j=v^*-c_0(x^*).
\]
This is a feasible bounded linear system of polynomial dimension over
$F$. Its coefficients have polynomial descriptions, because evaluating
the input polynomials uses bounded degree and the sampled coordinates
have polynomial descriptions in the same field. The algebraic linear
programming result of \citet[Section~5, Remark~1]{s4:adler1994},
which outlines the rational-machine implementation of their algorithm,
solves this system in time polynomial in the LP dimension, coefficient
encoding, and degree of a common extension. That degree is polynomial
here; separate small representations of unrelated algebraic
coefficients would not suffice. The explicit coordinate-box rows ensure
full column rank of the recovery constraint matrix. All recovered
blocks may be chosen in $F$, because a vertex is obtained from a
nonsingular subsystem by field linear algebra. Determinant and
field-arithmetic bounds give polynomial encoding length, consistent
with the algorithm's output guarantee. Explicitly, represent each
coefficient in the power basis of $F$ and clear its rational
denominators. A determinant of order at most the recovery dimension
has polynomial degree and coefficient bit length before reduction modulo
the minimal polynomial of $\alpha$: its factorial number of expansion
terms contributes only polynomially many bits. Polynomial division and
inversion in $F$ preserve polynomial encoding at these dimensions.
This bounds the Cramer ratios for an entire vertex, rather than
assuming that an arbitrary sequence of field operations has small
intermediate bit length. Together with $(x^*,v^*)$ these
blocks give the promised optimizer.

\emph{Aggregate inequalities.}
A row $\sum_j a_j(x)^Tz_j\le\beta(x)$ is converted to an equation by
adding its nonnegative slack as a scalar block. A polynomial-bit
rational upper bound for that slack can be computed, even if the core
formula does not explicitly give a box. First decide whether $C$ is
empty. If nonempty, eliminate all but one coordinate from its formula,
for each of the fixed number of coordinates. Each compact coordinate
projection has finite endpoints, and each endpoint is a root of a
nonzero polynomial in the resulting univariate formula: on an interval
containing no such root every atomic sign, and hence membership, is
constant. A Cauchy root bound for these polynomially encoded
polynomials gives a rational box $[-R,R]^r$ containing $C$, with
polynomial bit length. When $r=0$ no such step is needed. Interval
arithmetic over this box and the supplied block boxes bounds
$\beta(x)-\sum_j a_j(x)^Tz_j$ by a polynomial-bit rational number in
absolute value. The coefficient degree is fixed, and the polynomial
number of summands adds only polynomial encoding. A fixed number of
such slack blocks preserves all dimension hypotheses. This completes
the proof.
\end{proof}

The block LP algorithms of
\citet[Theorem~1]{s4:cslovjecsek2021} also exploit a fixed number of
linking rows and independent local optimization. Their matrices and
polyhedra are fixed linear data. The additional step above is global
optimization over a nonlinear core that changes those local data;
fixing a few nonlinear variables in an otherwise unrestricted LP
family would not establish this theorem.

\begin{remark}[What the theorem requires]\label{s4:core-limits}
A fixed number of design binaries may be included in the core by
$x_i(x_i-1)=0$, provided the rest of the model retains the displayed
form. The theorem does not cover an unbounded family of local binaries:
with no core, blocks $z_j\in\{0,1\}$ and one aggregate equation already
encode subset sum. Nor can arbitrary compact convex nonlinear blocks
replace the polyhedra without changing the output guarantee. Blocks
$0\le z_j\le a_j+1$, $z_j^2\le a_j$, with objective $-\sum_jz_j$,
have optimum $-\sum_j\sqrt{a_j}$. For distinct primes $a_j$ this value
has degree $2^N$. Indeed the multiquadratic extension has independent
sign automorphisms (the prime square classes are independent), and
its distinct square-root monomials are a rational basis. No nontrivial
sign change fixes the displayed sum, so its orbit has size $2^N$.
This is an exact-output obstruction, not an approximation hardness
claim. Polyhedral support values avoid it because they are rational
functions of one common fixed-dimensional core. The theorem also gives
no strongly polynomial or fixed-parameter runtime bound.
\end{remark}

\subsection{Two basic pooling consequences}\label{s4:pooling-basic}

\begin{corollary}[Fixed numbers of inputs and pools]\label{s4:fixed-inputs}
For fixed $m,p$, pooling with at most $m$ inputs and $p$ pools admits
polynomial-bit exact global optimization and optimizer recovery.
The numbers of products and qualities, and the bypass arcs, are
unrestricted.
\end{corollary}
\begin{proof}
If there are no inputs, check the all-zero flow directly. Otherwise,
remove a pool with no incoming or no outgoing arc after forcing its
incident flows to zero and checking that zero belongs to its throughput
interval and every incident-arc interval.
Keep every adjacent external-node contract. For each remaining pool
introduce inlet fractions $\theta_{i\ell}\ge0$ with
$\sum_i\theta_{i\ell}=1$ and $\theta_{i\ell}=0$ on missing inlet arcs.
They form a compact core of dimension at most $mp$.

The block for product $j$ contains $(v_{\ell j})_\ell$ and
$(z_{ij})_i$, of dimension at most $p+m$, with missing arcs fixed to
zero. Its flow boxes and product-throughput bounds are linear. For
attribute $k$, the upper quality row is
\begin{equation}
 \sum_\ell\left(\sum_i\lambda_{ik}\theta_{i\ell}
                    -\overline\mu_{jk}\right)v_{\ell j}
 +\sum_i(\lambda_{ik}-\overline\mu_{jk})z_{ij}\le0.
 \label{s4:fraction-product}
\end{equation}
The lower row replaces these coefficients by their lower-bound
counterparts $\underline\mu_{jk}-\sum_i\lambda_{ik}\theta_{i\ell}$
and $\underline\mu_{jk}-\lambda_{ik}$. Thus an arbitrary attribute
count increases the number of local rows only.

Let $T_\ell=\sum_jv_{\ell j}$. The remaining interval bounds apply to
\begin{equation}
 T_\ell,\qquad
 y_{i\ell}:=\sum_j\theta_{i\ell}v_{\ell j},\qquad
 s_i=\sum_j\left(\sum_\ell\theta_{i\ell}v_{\ell j}+z_{ij}\right).
 \label{s4:fraction-aggregates}
\end{equation}
These impose pool capacities, inlet-arc bounds, and input withdrawal
bounds, respectively, using at most $2(p+mp+m)$ aggregate inequalities.
A feed-arc cost $c_{i\ell}y_{i\ell}$ contributes
$\sum_i c_{i\ell}\theta_{i\ell}$ to the coefficient of every
$v_{\ell j}$; all other costs already belong to their product blocks.
Theorem~\ref{s4:fixed-core} applies with affine coefficient data.

For exactness, a positive-throughput physical pool has fractions
$y_{i\ell}/T_\ell$. A zero-throughput pool permits any simplex vector
on its nonempty incoming-arc set. Conversely,
$y_{i\ell}=\theta_{i\ell}T_\ell$ gives mass balance and a pool quality
$\sum_i\lambda_i\theta_{i\ell}$. Equations
\eqref{s4:fraction-product}--\eqref{s4:fraction-aggregates} are then
exactly the physical quality and throughput conditions, at positive
and zero throughput. The objective is preserved in both directions.
\end{proof}

\begin{corollary}[Fixed pools and attributes without bypasses]
\label{s4:fixed-qualities}
For fixed $p,K$, pooling with at most $p$ pools and $K$ attributes,
without bypasses, admits polynomial-bit exact global optimization and
optimizer recovery. Inputs and products may be arbitrarily numerous.
A fixed number of bypass arcs is also permitted.
\end{corollary}
\begin{proof}
If there are no inputs, all flows vanish and feasibility is checked
directly. Otherwise the $pK$ pool qualities form the core, bounded
coordinatewise between the smallest and largest input qualities.
For each input, use the block $(y_{i\ell})_\ell$ with its arc and
withdrawal bounds. For each product, use the block $(v_{\ell j})_\ell$
with its arc, delivery, and quality bounds. All local dimensions are
at most $p$, and each product quality row is affine in the core.
The only linking equations are
\begin{equation}
 \sum_i y_{i\ell}-\sum_jv_{\ell j}=0,\qquad
 \sum_i\lambda_{ik}y_{i\ell}-q_{\ell k}\sum_jv_{\ell j}=0.
 \label{s4:no-bypass-balances}
\end{equation}
There are $p(K+1)$ equations; pool throughput bounds add at most $2p$
aggregate inequalities. Costs separate over the blocks. These are the
physical concentration equations, with valid arbitrary boxed qualities
at inactive pools. Theorem~\ref{s4:fixed-core} proves the claim.

For a fixed number of bypasses, put their bounded flows in the core.
Subtract their input withdrawals, product deliveries, and known quality
masses from the appropriate local right-hand sides, and put their
costs in the core objective. The dimensions remain fixed and the
reformulation remains exact. Unbounded bypasses require a further
structural argument, supplied next.
\end{proof}

The piecewise-parametric pooling analysis of
\citet[Assumption~2.2]{s3:baltean2018} is a related decomposition
method. It uses one quality, removes feed and pool capacities, and
fixes positive product demands, unlike the capacitated,
arbitrary-attribute model of Corollary~\ref{s4:fixed-inputs}.

These corollaries extend the one-pool, fixed-input algorithms of
\citet[Theorem~1]{s4:boland2017} and
\citet[Theorem~4.1]{s3:haugland-hendrix2016}, whose models omit
bypasses; the latter discusses this convention explicitly. For its
final no-bypass model, \citet[Section~6, p.~214]
{haugland2016-the-computational-complexity-of-the} asks about a fixed
pool count together with a bound on inputs, products, or qualities.
Corollaries~\ref{s4:fixed-inputs} and~\ref{s4:fixed-qualities} answer
the input-bounded and quality-bounded cases, respectively. For the
remaining case, in which only pools and products are bounded and the
numbers of inputs and qualities grow, Theorem~\ref{s3:two-pools} gives
ordinary NP-completeness with two pools and two products, without
bypasses. The two-source/two-product hardness results in the same
Haugland article allow the pool count to grow, so they do not
contradict the fixed-input conclusion here. The earlier one-pool
algorithms and the fixed-dimensional Minkowski and algebraic algorithms
are prior ingredients; the uniform block theorem and its pooling
mappings are the additional results proved here.

\subsection{Affine quality rank and bypass vertex integrity}
\label{s4:bypass-structure}

The preceding block decompositions fail for unrestricted bypass graphs,
because a bypass can couple an input block to a product block. A small
set of bypass endpoints can absorb these couplings if the remaining
components are small. Low affine rank of the input qualities lets us
retain all product specifications with only a fixed number of aggregate
masses.

For a graph $G$, define its vertex integrity by
\[
 \iota(G)=\min_{W\subseteq V(G)}
 \left(|W|+\max_{H\text{ component of }G-W}|V(H)|\right),
\]
where the maximum over no components is zero. We use this established
graph parameter only through its definition. For the polynomial-time
results with fixed parameters here, a bounded value is equivalent to fixed
integers $c,h$ such that deleting at most $c$ vertices leaves components
of at most $h$ vertices. A vertex cover of size $c$ gives $h=1$.

\begin{theorem}[Controlled bypass structure]\label{s4:vertex-integrity}
Fix the pool count $p$, the affine rank $t$ of the rational input-quality
vectors, and a bound on the vertex integrity of the bypass graph.
Pooling feasibility and global optimization then admit polynomial-bit
exact algorithms, with a polynomial-size algebraic optimizer when
feasible. Inputs, products, attributes, and bypass arcs may all grow.
The product specification vectors are unrestricted.
\end{theorem}
\begin{proof}
The empty-input instance has zero flow and is checked directly. Otherwise
use the exact affine coordinates of Section~\ref{s1:encoding}:
\begin{equation}
 \lambda_i=\lambda_0+Ba_i,\qquad q_\ell=\lambda_0+B\eta_\ell,
 \qquad B\in\Q^{K\times t}.
 \label{s4:quality-coordinates}
\end{equation}
The columns of $B$ are independent and all rational encoding lengths
are polynomial. Put
$\min_i a_{is}\le\eta_{\ell s}\le\max_i a_{is}$ for every coordinate.
Active pools have these averaged coordinates; inactive pools may be
assigned arbitrary coordinates in the same box. The pool equations
are mass balance and
$\sum_i a_{is}y_{i\ell}=\eta_{\ell s}\sum_jv_{\ell j}$.
Their equivalence to all original attribute balances follows by
multiplying by $B$ and using its column independence.

Choose deleted inputs $A\subseteq I$ and products $B_0\subseteq J$,
with $c_I=|A|$, $c_J=|B_0|$ and $c_I+c_J\le c$, leaving components
$(I_\beta,J_\beta)$ of at most $h$ vertices. This can be found by
enumerating all vertex subsets of size at most $c$ and checking their
components, in polynomial time for fixed $c$. Failure of this search
means only that the structural hypothesis is absent, not that the
pooling instance is infeasible. Isolated vertices are components.

The core contains all $pt$ coordinates $\eta_{\ell s}$, all existing
$y_{i\ell}$ with $i\in A$, all $v_{\ell j}$ with $j\in B_0$, and all
$z_{ij}$ with $(i,j)\in A\times B_0$. It also contains, for each
$j\in B_0$, a total delivery variable $\tau_j$ and coordinate-mass
variables $M_{js}$, $1\le s\le t$. Each component block owns
\begin{equation}
 \begin{gathered}
 y_{i\ell}\ (i\in I_\beta),\qquad
 v_{\ell j}\ (j\in J_\beta),\\
 z_{ij}\ ((i,j)\in I_\beta\times J_\beta),\qquad
 z_{ij}\ ((i,j)\in A\times J_\beta),\qquad
 z_{ij}\ ((i,j)\in I_\beta\times B_0),
 \end{gathered}\label{s4:ownership}
\end{equation}
restricted to existing arcs. Every flow is assigned exactly once: no
bypass joins two distinct remaining components. Each block has at
most $ph+ch+h^2$ coordinates. For $h=1$, the sharper bound is $p+c$.

All arc bounds belong to the owning core or block. For each input in
$I_\beta$, all its incident flows are in block $\beta$, so its withdrawal
bounds are local. For each product in $J_\beta$, the same is true of
its delivery and attribute rows. For example, its upper row is
\begin{equation}
 \sum_\ell(\lambda_{0k}+(B\eta_\ell)_k-\overline\mu_{jk})v_{\ell j}
 +\sum_i(\lambda_{ik}-\overline\mu_{jk})z_{ij}\le0.
 \label{s4:local-quality}
\end{equation}
The lower row reverses the corresponding quality differences. Given
the core these are linear inequalities in the block coordinates.
An arbitrary attribute count creates only more local rows.

The remaining pool mass and coordinate equations are
\begin{align}
 \sum_\beta\left(\sum_{i\in I_\beta}y_{i\ell}
                  -\sum_{j\in J_\beta}v_{\ell j}\right)
 &=\sum_{j\in B_0}v_{\ell j}-\sum_{i\in A}y_{i\ell},\\
 \sum_\beta\left(\sum_{i\in I_\beta}a_{is}y_{i\ell}
             -\eta_{\ell s}\sum_{j\in J_\beta}v_{\ell j}\right)
 &=\eta_{\ell s}\sum_{j\in B_0}v_{\ell j}
                  -\sum_{i\in A}a_{is}y_{i\ell}.
 \label{s4:pool-aggregates}
\end{align}
There are $p(t+1)$ equations, with affine block coefficients and
core right-hand sides of degree at most two. Pool bounds apply to
$\sum_{j\in B_0}v_{\ell j}+\sum_\beta\sum_{j\in J_\beta}v_{\ell j}$.
For a deleted input $i\in A$, withdrawal bounds apply to
\[
 \sum_\ell y_{i\ell}+\sum_{j\in B_0}z_{ij}
                      +\sum_\beta\sum_{j\in J_\beta}z_{ij}.
\]
These contribute at most $2p+2c_I$ aggregate inequalities.

Deleted products need a separate step. Leaving all of their original
attribute rows as aggregate inequalities would make their number grow
with $K$, even when $t$ is fixed. Instead define their retained totals by
\begin{align}
 \sum_\beta\sum_{i\in I_\beta}z_{ij}
 &=\tau_j-\sum_\ell v_{\ell j}-\sum_{i\in A}z_{ij},\\
 \sum_\beta\sum_{i\in I_\beta}a_{is}z_{ij}
 &=M_{js}-\sum_\ell\eta_{\ell s}v_{\ell j}
                    -\sum_{i\in A}a_{is}z_{ij}.
 \label{s4:covered-totals}
\end{align}
These use $c_J(t+1)$ aggregate equations. All their delivery and quality
requirements now lie in the core: delivery bounds apply to $\tau_j$,
and the attribute rows are
\begin{equation}
 (\lambda_{0k}-\overline\mu_{jk})\tau_j+\sum_sB_{ks}M_{js}\le0,
 \qquad
 (\underline\mu_{jk}-\lambda_{0k})\tau_j-\sum_sB_{ks}M_{js}\le0.
 \label{s4:covered-quality}
\end{equation}
No rank hypothesis on the product specification vectors was used.
To bound these variables, let $H_j$ be the sum of the finite upper
bounds on all arcs entering $j$, and $A_s=\max_i|a_{is}|$. Valid bounds
are $0\le\tau_j\le H_j$ and $|M_{js}|\le A_sH_j$, because all incoming
coordinate qualities lie in $[-A_s,A_s]$. The bounds are rational of
polynomial encoding length. The core domain, formed from these boxes,
core arc bounds, and \eqref{s4:covered-quality} and delivery bounds,
is compact.

For exactness, map every physical flow to its uniquely owned
coordinates, choose its active pool averages, and set the retained
product totals to their actual values. All constructed equations hold.
Conversely, the local rows impose every undeleted external-node bound,
\eqref{s4:pool-aggregates} imposes the exact pool balances, and the
pool and deleted-input aggregate inequalities impose their remaining
throughput bounds. Equations \eqref{s4:covered-totals} identify the
actual deleted-product masses, so \eqref{s4:covered-quality} gives
exactly their original attribute requirements. At zero pool throughput
nonnegative incident flows vanish, and no division is needed. All
physical arc costs belong to their owning core or block objective and
are unchanged. Thus both formulations have the same feasible flows
and cost values.

The parameters of Theorem~\ref{s4:fixed-core}, after adding bounded
scalar slacks, satisfy
\begin{equation}
 \begin{aligned}
 r&\le pt+pc+c^2+c_J(t+1),\\
 d&\le\max(1,ph+ch+h^2),\\
 k&\le p(t+1)+c_J(t+1)+2p+2c_I,\qquad D=2.
 \end{aligned}\label{s4:structure-dimensions}
\end{equation}
All are fixed. If no blocks remain, append a zero block. The theorem
therefore supplies exact optimization and recovery; the rational map
\eqref{s4:quality-coordinates} recovers every original pool attribute
in polynomial bit time as well.
\end{proof}

\begin{corollary}[Specializations and elementary cases]
\label{s4:structural-corollaries}
Theorem~\ref{s4:vertex-integrity} includes each of the following cases:
\begin{enumerate}[label=(\roman*)]
 \item Fixed pools, fixed attribute count, and a fixed-size bypass
 vertex cover, allowing arbitrarily many bypasses at fixed input or
 product hubs.
 \item Fixed pools and fixed affine quality rank, with bypass components
 of bounded size, including an unbounded matching.
 \item Fixed pools and fixed affine quality rank without bypasses.
 A fixed number of distinct input-quality vectors also suffices in
 place of the rank bound.
 \item Fixed inputs and pools with arbitrary bypasses and attributes,
 as in Corollary~\ref{s4:fixed-inputs}.
\end{enumerate}
With no pools, or with affine quality rank zero, the model is an LP
without any bypass-structure restriction; in the rank-zero case the
pool count can also grow.
\end{corollary}
\begin{proof}
For (i), $t\le K$ and a vertex cover leaves isolated vertices. For (ii)
take $c=0$. For (iii), a graph with no edges has singleton components,
and $q$ distinct vectors have affine rank at most $q-1$. For (iv), all
inputs form a vertex cover and $t\le|I|-1$. With no pools only known
input-quality bypasses remain. At rank zero all input vectors equal
$\lambda_0$; set every pool quality to $\lambda_0$. Conservation then
makes all mixing equations redundant and all product conditions linear.
Inactive pools cause no exception.
\end{proof}

The covered-product totals in \eqref{s4:covered-totals} are what permit
affine rank in place of the attribute count: replacing the attribute
count by rank in a row-count bound alone would overlook the arbitrary
specification rows at a covered product. Graph degree and treewidth
alone do not bound vertex integrity, since paths already have growing
vertex integrity. The next subsections treat path geometry directly,
with a different retained boundary and a more restricted objective.

\subsection{Composition of planar relations}\label{s4:planar}

A bypass path can be described by its consecutive scalar arc flows.
Each internal node gives a compact planar relation between its two
incident flows. We show that projecting a whole path onto its two
endpoint flows has polynomial description size. This projection
statement is stronger than deciding a fixed linear instance, and it
retains no accumulated objective coordinate.

\begin{lemma}[Two planar relations]\label{s4:polygon-composition}
Let $P\subset\R^2$ in coordinates $(x,y)$ and $Q\subset\R^2$ in
coordinates $(y,z)$ be compact rational polyhedra, given by $m_P$ and
$m_Q$ weak linear inequalities. Their composition
\[
 P\circ Q=\{(x,z):\exists y,\ (x,y)\in P,\ (y,z)\in Q\}
\]
has a rational inequality description with at most
$64(m_P+m_Q+2)$ rows. It can be constructed with polynomially many
rational operations. If the input row coefficients have bit lengths
at most $b_P,b_Q$, the output coefficients can be chosen with bit
length at most $C(b_P+b_Q+1)$ for a universal constant $C$.
Empty sets, segments, and points are included.
\end{lemma}
\begin{proof}
First compute the interval $J=\operatorname{proj}_y Q$ and clip $P$
by $y\in J$. This adds at most two rows. Empty cases are immediate.
Suppose initially that the resulting $x$ domain and $J$ are
nondegenerate intervals. Write the vertical slices of clipped $P$ as
$[l(x),u(x)]$ and those of $Q$ as $[q(y),r(y)]$. The functions $l,q$
are convex piecewise affine, and $u,r$ are concave piecewise affine.
Each of $l,u$ has at most $m_P+2$ affine pieces, and each of $q,r$
has at most $m_Q$ pieces: each piece comes from a nonvertical
supporting row of the relevant polygon. These statements also hold
when the lower and upper boundaries coincide.

The minimum attainable $z$ for a given $x$ is
$g(x)=\min_{l(x)\le y\le u(x)}q(y)$. Let $[a,b]\subseteq J$ be the
minimizer interval of $q$, and $q_{\min}$ its minimum. Define on $J$
\[
 D(t)=\begin{cases}q(t)&t<a,\\q_{\min}&t\ge a,\end{cases}
 \qquad
 L(t)=\begin{cases}q_{\min}&t\le b,\\q(t)&t>b.\end{cases}
\]
The function $D$ is convex and nonincreasing; $L$ is convex and
nondecreasing. The interval $[l(x),u(x)]$ either lies to the left of
$[a,b]$, meets it, or lies to its right. In the three cases its
minimum of $q$ is respectively $q(u(x))$, $q_{\min}$, or $q(l(x))$.
Equivalently,
\begin{equation}
 g(x)=\max\{L(l(x)),D(u(x))\}.
 \label{s4:clamped-envelopes}
\end{equation}
Both terms are convex: a convex nondecreasing outer function preserves
convexity, and a convex nonincreasing outer function composed with a
concave inner function is convex. These facts follow directly from
the defining convexity inequality and the indicated monotonicity.

To count pieces of a composition, subdivide the inner domain at its
own knots and at the boundary points of the preimages of every outer
knot. A level set of a univariate convex or concave function has at
most two boundary points: away from a minimum or maximum plateau
there are at most two crossings, while a plateau contributes only
its two endpoints. Thus, if the inner and outer functions have $n$
and $m$ affine pieces, this subdivision has at most $n+2m+2$ pieces,
using a loose bound that also covers domain endpoints. On each open
piece the composition is affine; flat preimages introduce no extra
formula. Each of $D,L$ has at most $m_Q+1$ pieces. Consequently each
term in \eqref{s4:clamped-envelopes} has at most
$m_P+2m_Q+6$ pieces.

A convex piecewise-affine function on an interval equals the maximum
of the affine extensions of its pieces: each such line supports the
whole function. Therefore the maximum of two such functions uses at
most the sum of their numbers of pieces. It follows that $g$ has at
most $2m_P+4m_Q+12$ pieces. Reflecting $z$ gives the same bound for
the concave upper boundary. The feasible $z$ at fixed $x$ form an
interval, since they are the projection of a nonempty convex slice.
Hence these two envelopes, with two vertical domain bounds, describe
the composition. The resulting row count is at most
$4m_P+8m_Q+26$, below the stated loose bound.

For degeneracies, if the $x$ domain is a singleton, the composition
is a vertical interval or point; minimize $q$ and maximize $r$ over
the corresponding feasible $y$ interval. If $J$ is a singleton,
clipping fixes $y$ and the composition is the Cartesian product of
an $x$ interval and a $z$ interval. These need at most four rows.
A nonvertical segment has coincident affine lower and upper envelopes
and was already covered. The empty relation can be described by
$0\le-1$.

For a constructive algorithm and its coefficient bound, combine the
row descriptions of $P(x,y)$ and $Q(y,z)$ and perform one
Fourier--Motzkin elimination of $y$. Keep rows with zero coefficient
on $y$; combine every pair with opposite coefficient signs after
positive scaling to cancel $y$. This gives at most
$(m_P+m_Q)^2+m_P+m_Q$ rows in $(x,z)$ and is an exact projection.
Each new row is formed from two input rows by a fixed number of
rational additions and multiplications. Pairs may come from the
same child as well as different children; either case satisfies the
stated $C(b_P+b_Q+1)$ bit bound.

Compute the bounded planar result's vertices by pairwise line
intersections and row feasibility checks. A nonempty bounded
polyhedron has vertices, even if its affine dimension is smaller.
In full dimension, remove redundant rows one at a time, retaining one
row for each facet; the envelope argument bounds their number. Testing
all rows against the original list and deleting them simultaneously
would be invalid in the presence of duplicate rows. Redundancy is
decidable by optimizing each candidate row over the system with that
row removed, using an exact planar LP (which may be unbounded). In dimension one, choose the two extreme endpoints from
the vertices and describe their segment by a line equation and two
endpoint inequalities; in dimension zero, fix the two coordinates.
These exceptional descriptions also require only a constant number
of rational operations on input-row intersections per coefficient.
All cases therefore give the claimed output size, polynomial
arithmetic work, and universal coefficient-bit bound.
\end{proof}

\begin{theorem}[Exact path endpoint projection and lifting]
\label{s4:path-projection}
For a chain of compact rational planar relations
\begin{equation}
 (x_{i-1},x_i)\in P_i\qquad(1\le i\le n),
 \label{s4:path-chain}
\end{equation}
the exact projection onto $(x_0,x_n)$ has a rational inequality
description of polynomial size and bit length, computable in
polynomial bit time. A feasible endpoint pair can be lifted to all
internal coordinates using polynomially many arithmetic operations and
comparisons. If the endpoints lie in a common represented real-algebraic
field of polynomial degree and
encoding length, the lifted coordinates stay in that field and have
polynomial total encoding length; the lift takes polynomial bit time,
including the input and output encoding.
\end{theorem}
\begin{proof}
Compose consecutive subchains in a balanced binary tree. Write
$s(P)=m_P+2$ for a relation's row count plus two. By
Lemma~\ref{s4:polygon-composition}, the parent size is at most a
universal constant times the sum of its children's sizes. The depth
is at most $\lceil\log_2 n\rceil$. At depth $a$ the total size is
at most a constant to the power $a$ times the total leaf size.
Therefore the maximum and total intermediate sizes are
$n^{O(1)}\sum_i(m_i+2)$. Each composition takes polynomial time in
its children's sizes, so the total arithmetic work is polynomial.
The same recurrence applies to rational coefficient bit length:
$b_{\rm parent}\le C(b_{\rm left}+b_{\rm right}+1)$. Logarithmic
depth makes the resulting bit length polynomial in the original
encoding. Rational comparisons, line intersections, and redundancy
checks consequently all have polynomial bit cost. This argument
uses balanced composition; no comparable bound for an arbitrary
sequential elimination order is needed.

Store every child relation in the tree. At an internal node with
feasible endpoints $(x,z)$, let $Y_L(x)$ be the left child's slice
of shared $y$ values and $Y_R(z)$ the right child's slice. These
are closed bounded intervals, and their intersection is nonempty
by the definition of composition. Choose the midpoint of that
intersection and recurse into both children. This proves exact
lifting, including zero-length slices.

For the representation guarantee, each endpoint of a slice is a
maximum or minimum of rational affine functions of the prescribed
outer coordinate. Rows whose coefficient on $y$ is zero are
feasibility tests; every division used to compute a finite slice
bound is by a nonzero rational row coefficient. Exact comparisons
select the active interval bounds. Their midpoint is a rational
affine combination of the two parent endpoints. Recursing along
the balanced tree therefore expresses every recovered coordinate
as a rational affine combination of the original two endpoints,
with polynomial coefficient length by the same depth bound. No
new algebraic root or field extension is required. Exact sign
comparisons in the supplied common field have polynomial bit
complexity in its degree and representation size, proving the last
claim.
\end{proof}

Pairwise projection by resultants is established in the literature on
two variables per inequality; see the revised manuscript of
\citet[Section~3]{s4:simon2010}. The envelope and elimination methods
of \citet[Section~2]{s4:hochbaum1994} are further algorithmic
antecedents. The path proof above does not rely on a general
polynomial-size closure claim for all such systems; it controls both
the number of rows and the bit length for the two-endpoint projection
that pooling needs.

\subsection{Bounded pool attachments and arbitrary attributes}
\label{s4:attachments}

A bounded number of pool-incident arcs gives a fixed physical interface
with the bypass graph, even if its path components are long. The input
qualities on those paths are known, so their internal relations are
rational and independent of the pool mixture.

\begin{theorem}[Bounded attachments through degree-two bypasses]
\label{s4:bounded-attachments}
Fix the total number $a$ of input--pool and pool--product arcs.
If the bypass graph has maximum degree at most two, pooling feasibility
is decidable in polynomial bit time, with exact algebraic feasible-flow
recovery. The number of attributes, distinct input-quality vectors,
inputs, products, and bypass arcs is unrestricted. Multiple pools and
all the rational throughput, arc, and quality bounds specified at the
start of this section are allowed.
\end{theorem}
\begin{proof}
A pool with no incoming or no outgoing arc has zero throughput and
zero incident flows. Check that zero belongs to its throughput interval
and every incident-arc interval, then remove these arcs and the pool.
Preserve each adjacent input or product requirement, since other arcs may satisfy
it. If no pool arcs remain, the instance is a rational blending LP.
We may thus assume every remaining pool has an inlet and an outlet.

Call an external node adjacent to a pool an attachment. There are
at most $a$ distinct attachments, including shared nodes adjacent to
several pools. Retain in the core all at most $a$ pool-incident flows,
all bypass flows incident to any attachment, and a fraction
$\theta_{i\ell}$ for every existing inlet arc. At most $2a$ distinct
bypass arcs are retained because each attachment has bypass degree
at most two; an arc between two attachments is retained only once.
There are at most $a$ fractions. Thus the core has at most $4a$
scalar coordinates. Put all flow bounds in the core for retained
arcs and require
\begin{equation}
 \theta_{i\ell}\ge0,\qquad
 \sum_{i:i\ell\in A}\theta_{i\ell}=1,\qquad
 \sum_i y_{i\ell}=\sum_jv_{\ell j},\qquad
 y_{i\ell}=\theta_{i\ell}\sum_{i'}y_{i'\ell}.
 \label{s4:attachment-fractions}
\end{equation}
Here the arc set $A$ has its original meaning from Section~\ref{s1:model}.
For each pool define its full quality vector by
$q_\ell=\sum_i\lambda_i\theta_{i\ell}$; this is an expression, not
$K$ new variables. All input or product throughput rows at attachments
involve retained flows only. So do their quality rows, after this
substitution, and all pool throughput bounds. Their degrees are at
most two, even when $K$ grows. At positive throughput the equations
force the actual mixture fractions. At zero throughput any legal
simplex fraction gives the same zero transmitted mass.

Delete the attachment nodes from the bypass graph. Every remaining
component is a path, a cycle, or an isolated node. A component incident
to a deleted node is a path, possibly an isolated node, with at most
two boundary arcs: a remaining vertex on a cycle already has its full
two neighbors within that cycle, so such a cycle cannot meet the
boundary. For a path, only its endpoints have room for a boundary
arc. The flows on these boundary arcs have been retained in the core.
Every nonattachment input has only bypass flows and hence has a
withdrawal row in its at most two adjacent scalar arc coordinates.
Every nonattachment product likewise has only known-quality bypass
flows; its delivery and each attribute row involve its at most two
adjacent coordinates, with rational constant coefficients. Together
with the finite arc boxes, each node thus gives a compact planar
relation between successive flows along the path.

For a component with two boundary arcs, apply
Theorem~\ref{s4:path-projection} and insert its exact endpoint
inequalities into the core. A component with one boundary arc has
an interval projection; append a dummy endpoint fixed to zero at its
free end and express the terminal node's univariate rows as a
planar relation with this dummy coordinate. The same theorem gives
the required projection after substituting zero. Components with no
boundary, including detached cycles, are rational LP feasibility
problems and can be solved independently. A node with no usable arc
and a positive withdrawal or delivery lower bound is rejected by
these same local conditions. A bypass with both endpoints at
attachments is already wholly in the core. A bypass cycle cut by
an attachment can leave a path whose two boundary arcs meet that
same attachment; the two arc coordinates are still distinct scalar
endpoints, so the composition proof is unchanged.

The resulting core has fixed dimension, polynomially many
polynomially encoded rows, and degree at most two. Its feasible set
is exactly the projection of the bounded physical formulation with
simplex fractions: the forward implication restricts a physical
flow, and the converse lifts every projected path and combines it
with the independently feasible detached components. All attachment
rows and \eqref{s4:attachment-fractions} remain, so this converse
enforces every physical mixing, capacity, and quality requirement.
The core is compact, and fixed-dimensional real-algebraic decision
and sampling give a feasible core point or an infeasibility decision
in polynomial bit time.

Represent the sampled coordinates in one polynomial-degree field.
Lift all boundary-connected paths through their stored composition
trees using Theorem~\ref{s4:path-projection}; their rational affine
lifts preserve that field and polynomial encoding. Detached components
have polynomial-bit rational LP witnesses. Finally evaluate the affine
expressions for all pool qualities. This recovers an exact physical
feasible flow and proves the theorem.
\end{proof}

\begin{corollary}[Two feeds and two outlets]\label{s4:two-two-feasibility}
With one pool, at most two actual inlet arcs, at most two actual outlet
arcs, and bypass degree at most two, pooling feasibility and exact
feasible-flow recovery are polynomial-time problems, even with
arbitrarily many attributes.
\end{corollary}
\begin{proof}
Theorem~\ref{s4:bounded-attachments} applies with $a\le4$. A sharper
core count is thirteen: four pool flows, at most eight incident bypass
flows, and one inlet fraction $\theta\in[0,1]$. With two feeds set
$q=\theta\lambda_1+(1-\theta)\lambda_2$ and
$y_1=\theta(y_1+y_2)$, together with mass balance. With one feed the
fraction is fixed and the core is smaller. All zero-throughput and
missing-arc cases are covered by the preceding proof.
\end{proof}

\subsection{Objectives on a fixed set of retained coordinates}
\label{s4:bounded-support}

\begin{theorem}[Bounded-support exact optimization]
\label{s4:fixed-support}
Under Theorem~\ref{s4:bounded-attachments}, fix also an integer $s$.
Designate at most $s$ further actual arc-flow coordinates. Any rational
polynomial objective of fixed degree in these coordinates and the
retained attachment core can be globally optimized in polynomial bit
time, with exact optimum and physical optimizer recovery. Polynomially
many additional closed polynomial equalities and weak inequalities
of fixed degree on those same retained coordinates are allowed.
In particular, the conclusion covers a linear cost supported on a
fixed number of actual arcs.
\end{theorem}
\begin{proof}
Retain every designated flow in the core and cut the scalar relation
chains at it. A formerly eliminated path then splits into segments
with at most two retained endpoints. In a detached cycle, cutting at
one designated arc gives a path relation $R(x,x)$: the two endpoint
occurrences refer to the same retained flow. Several designated arcs
split the cycle into ordinary two-endpoint segments. Isolated nodes,
one-endpoint segments, and components without retained coordinates
are handled as in Theorem~\ref{s4:bounded-attachments}. This cuts
relations at coordinates, so local node constraints remain attached
to their correct pairs; none is discarded.

The projected system has at most $4a+s$ variables, polynomial
encoding, and fixed degree, including the new constraints. It is the
exact projection of a closed bounded feasible set, and hence compact.
Its continuous polynomial objective attains an optimum if it is
feasible. Introduce an objective-value variable, eliminate quantifiers
to describe attainable values, and sample an optimal fiber; equivalently
use the optimality formula \eqref{s4:optimal-sample} with this
fixed-dimensional system. Fixed-dimensional real-algebraic algorithms
have polynomial bit complexity and return a polynomial-size common-field
sample. Affine path lifting recovers all eliminated flows while
preserving the objective, because every coordinate on which it
depends was retained. Strict additional constraints are excluded from
this attainment statement.
\end{proof}

\begin{corollary}[An objective reducible by exact equalities]
\label{s4:equality-support}
Let $Ef=e$ be rational linear conservation or contract equations
valid for every feasible flow in an instance of
Theorem~\ref{s4:bounded-attachments}. Suppose a supplied rational
identity
\begin{equation}
 c=E^T\gamma+d\label{s4:objective-reduction}
\end{equation}
has polynomial encoding length, where $d$ is supported on a fixed
number of arcs. Then exact minimization of $c^Tf$ is polynomial-time
solvable, with value $\gamma^Te+\min d^Tf$ and physical optimizer
recovery.
\end{corollary}
\begin{proof}
Check \eqref{s4:objective-reduction} exactly by rational arithmetic.
For every feasible $f$, $c^Tf=\gamma^Te+d^Tf$; apply
Theorem~\ref{s4:fixed-support} to $d$. The equality must hold throughout
the feasible set, not just at selected candidates or for tight
inequalities guessed to hold at an optimum.
\end{proof}

For example, along a bypass path whose consecutive inputs have exact
supplies $S_i$, write their two arcs to consecutive products as
$S_i-u_i,u_i$. With product revenues $R_{i-1},R_i$, the contribution
of these arcs to revenue is
\begin{equation}
 \sum_i\bigl(R_{i-1}(S_i-u_i)+R_i u_i\bigr)
 =\sum_iR_{i-1}S_i+\sum_i(R_i-R_{i-1})u_i.
 \label{s4:price-changes}
\end{equation}
Their input production costs are constant under the exact supplies.
Thus a fixed total number of revenue changes along these paths gives
bounded support, with any attachment contributions already retained.
The number of distinct revenues alone does not bound the number of
changes: two values may alternate along an arbitrarily long path.

\subsection{A sharp feasibility boundary and the remaining scope}
\label{s4:scope}

\begin{corollary}[Product degree two versus three]
\label{s4:degree-boundary}
Consider one-pool pooling with exactly two actual inlet arcs and two
actual outlet arcs, total input out-degree at most two, and positive
flow lower bounds permitted. If every product has total in-degree at
most two, feasibility is decidable in polynomial bit time for arbitrary
rational quality data. Allowing product total in-degree three makes
feasibility strongly NP-complete, even with one upper-bound scalar
quality, flow bounds in $\{0,1,2,3,4\}$, and quality data in
\eqref{s3:palette}.
\end{corollary}
\begin{proof}
For the positive branch the bypass graph is a subgraph of the full
external-node incidence structure, so its maximum degree is at most
two. Corollary~\ref{s4:two-two-feasibility} applies. It actually allows
some attachments of larger total degree when their bypass degree
remains at most two.

For the negative branch, take the physical construction in the proof
of Theorem~\ref{s3:constant-thm} before the final contract-completion
objective \eqref{s3:completion} is introduced. Retain its exact source
supplies, collector demands, and all other designated flow contracts.
The proof through \eqref{s3:physical-threshold} already shows that this
contracted network is feasible exactly when the positive-product source
instance is a yes instance. The source threshold is encoded by physical
binary linear circuits, not by an extra objective row that would
disappear when economics are removed.

The degree-two input and degree-three product bounds, the two actual
feeds and outlets, and the fixed quality alphabet all hold before
contract completion. Positive exact node contracts have values in
$\{1,2,3,4\}$, and all flow bounds lie in the stated alphabet.
The source's large rational data appear only in polynomial-size
circuit topology, so even unary encoding of the physical data has
polynomial size. This proves strong feasibility NP-hardness. With
one pool and one scalar quality, the fixed-dimensional linear-fiber
certificate in Corollary~\ref{s2:pooling-np} gives NP membership.
\end{proof}

The boundary concerns feasibility with prescribed nonzero requirements.
When every flow lower bound is zero, the zero flow is feasible, and a
profit threshold is a different decision problem. The positive result
remains valid when its data are restricted to the negative branch's
fixed alphabets.

The path algorithm retains a fixed number of actual flow coordinates.
A two-coordinate endpoint feasibility relation does not represent an
unrestricted linear cost or a dense additional aggregate over
eliminated arcs: interior flows with the same endpoints can have
different aggregate values, and a threshold on that value cannot be
checked after it is discarded. These theorems therefore do not resolve
general profit optimization for one pool with two feeds, two outlets,
and degree-two bypasses, and degree-two bypasses with unboundedly many
attachments remain open (Section~\ref{s6:open}).
